\section{Aplicación práctica: Stable Diffusion 3 y Flux}

En esta sección, el ponente ofrece una breve introducción sobre la aplicación de la tecnología Flow Matching en los modelos actuales más avanzados para generación de imágenes y video.

\subsection{Fundamentos de flujo de modelos para SD3 y Flux}

El ponente señala explícitamente que muchos de los modelos generativos de vanguardia actuales han \textbf{cambiado de} modelos de difusión a modelos de flujo:

\begin{itemize}
    \item \textbf{Stable Diffusion 3} (SD3): Aunque el nombre incluye ``diffusion'', en realidad utiliza entrenamiento basado en Flow Matching; la red predice la velocidad (velocity) en lugar del ruido.
    \item \textbf{Flux}: Modelo desarrollado por Black Forest Labs, que también se basa en la tecnología Rectified Flow.
\end{itemize}

\begin{importantbox}{Elecciones de diseño clave para SD3}
Las decisiones técnicas centrales de Stable Diffusion 3 incluyen:
\begin{itemize}
    \item \textbf{Entrenamiento con Rectified Flow}: Uso de mapeo de flujo condicional mediante interpolación lineal; la red genera una velocidad $v_\theta(\mathbf{x}_t, t)$.
    \item \textbf{Arquitectura MM-DiT}: Adopción de Multimodal Diffusion Transformer como red base.
    \item \textbf{Objetivo de trayectoria recta}: Búsqueda de trayectorias de generación más directas mediante técnicas como Reflow, reduciendo así el número de pasos de muestreo.
    \item \textbf{Predicción de velocidad (v-prediction)}: En lugar de predicción de ruido ($\epsilon$-prediction), lo cual ofrece una mejor estabilidad numérica en ciertos niveles de ruido.
\end{itemize}
\end{importantbox}

\subsection{Modelo Flux}

El ponente menciona que, desde la primera sesión del curso, se solicitó a los estudiantes generar imágenes utilizando el modelo Flux. Una variante clave de Flux es \textbf{Flux.1-schnell}, que emplea la tecnología Rectified Flow junto con Reflow para lograr generación de imágenes de alta calidad en un solo paso o con pocos pasos.

\begin{knowledgebox}{Modelos de consistencia vs.\ Rectified Flow: dos rutas tecnológicas para acelerar la generación}
En cuanto a la aceleración del muestreo en modelos de difusión/flujo, actualmente existen dos rutas competitivas principales:
\begin{itemize}
    \item \textbf{Consistency Model (Modelo de consistencia)}: Aprendizaje de un ``mapeo rápido'' sobre las trayectorias ODE mediante Consistency Distillation (CD) o Consistency Training (CT), permitiendo saltar directamente desde cualquier punto intermedio al final. Esto se discutió en la Lecture 8/9.
    \item \textbf{Rectified Flow + Reflow}: Reducción del número de pasos necesarios al hacer que las trayectorias mismas sean más rectas. Este es el tema central de esta sesión.
\end{itemize}

Los resultados experimentales comparativos muestran:
\begin{itemize}
    \item El Consistency Model, al generar en 1--2 pasos, alcanza un FID aproximado de 3--4 en ImageNet 64$\times$64.
    \item El 2-Rectified Flow también logra un FID similar en generación de 1--2 pasos.
\end{itemize}
Ambos métodos tienen sus propias ventajas y desventajas, y siguen siendo áreas de investigación activa.
\end{knowledgebox}

\subsection{Resumen de la sección}

El Flow Matching, especialmente las variantes de Rectified Flow, se ha convertido en la tecnología fundamental para los modelos generadores de imágenes más avanzados del momento (SD3, Flux). Su ventaja central radica en un entrenamiento sencillo, trayectorias controlables y la posibilidad de generar con pocos pasos.


